% ============================================================
%   Series de Tiempo — Documentación de Notebooks
%   Autor: Iñaki Castillo
%   Fecha: Junio 2026
% ============================================================

\documentclass[12pt, a4paper]{article}

% ── Codificación y fuentes ────────────────────────────────
\usepackage[utf8]{inputenc}
\usepackage[T1]{fontenc}
\usepackage[spanish]{babel}
\usepackage{lmodern}

% ── Geometría y diseño ───────────────────────────────────
\usepackage[top=2.5cm, bottom=2.5cm, left=3cm, right=3cm]{geometry}
\usepackage{setspace}
\setstretch{1.15}

% ── Matemáticas ──────────────────────────────────────────
\usepackage{amsmath, amssymb, amsthm}

% ── Código fuente ─────────────────────────────────────────
\usepackage{listings}
\usepackage{xcolor}

\definecolor{codebg}{HTML}{F7F7F7}
\definecolor{keywordcolor}{HTML}{0000CC}
\definecolor{commentcolor}{HTML}{008000}
\definecolor{stringcolor}{HTML}{BA2121}
\definecolor{rulecolor}{HTML}{CCCCCC}
\definecolor{headerbg}{HTML}{2C3E50}

\lstdefinestyle{python}{
  language=Python,
  backgroundcolor=\color{codebg},
  basicstyle=\ttfamily\small,
  keywordstyle=\color{keywordcolor}\bfseries,
  commentstyle=\color{commentcolor}\itshape,
  stringstyle=\color{stringcolor},
  breaklines=true,
  breakatwhitespace=true,
  showstringspaces=false,
  frame=single,
  rulecolor=\color{rulecolor},
  xleftmargin=1em,
  xrightmargin=1em,
  numbers=left,
  numberstyle=\tiny\color{gray},
  numbersep=6pt,
  tabsize=4,
  captionpos=b
}
\lstset{style=python}

% ── Gráficos y colores ───────────────────────────────────
\usepackage{graphicx}
\usepackage{tikz}
\usepackage{tcolorbox}
\tcbuselibrary{skins, breakable}

% Caja de definición
\newtcolorbox{definicion}[1][]{
  enhanced, breakable,
  colback=blue!5!white,
  colframe=blue!50!black,
  fonttitle=\bfseries,
  title=Definición,
  #1
}

% Caja de nota/observación
\newtcolorbox{nota}[1][]{
  enhanced, breakable,
  colback=orange!8!white,
  colframe=orange!70!black,
  fonttitle=\bfseries,
  title=Nota,
  #1
}

% Caja de resultado
\newtcolorbox{resultado}[1][]{
  enhanced, breakable,
  colback=green!5!white,
  colframe=green!50!black,
  fonttitle=\bfseries,
  title=Resultado clave,
  #1
}

% ── Hipervínculos ─────────────────────────────────────────
\usepackage{hyperref}
\hypersetup{
  colorlinks=true,
  linkcolor=blue!60!black,
  urlcolor=blue!70!black,
  citecolor=red!60!black,
  pdftitle={Series de Tiempo — Iñaki Castillo},
  pdfauthor={Iñaki Castillo}
}

% ── Encabezados y pies ───────────────────────────────────
\usepackage{fancyhdr}
\pagestyle{fancy}
\fancyhf{}
\fancyhead[L]{\small Series de Tiempo}
\fancyhead[R]{\small Iñaki Castillo}
\fancyfoot[C]{\thepage}
\renewcommand{\headrulewidth}{0.4pt}

% ── Tablas y figuras ─────────────────────────────────────
\usepackage{booktabs}
\usepackage{array}
\usepackage{caption}

% ── Utilidades ───────────────────────────────────────────
\usepackage{enumitem}
\usepackage{parskip}

% ─────────────────────────────────────────────────────────
\begin{document}

% ── Portada ───────────────────────────────────────────────
\begin{titlepage}
  \centering
  \vspace*{2cm}
  {\Huge\bfseries Series de Tiempo:\\[0.3em]
   \Large Modelado, Análisis y Pronóstico\par}
  \vspace{1.5cm}
  {\large Documentación de Notebooks de Python\par}
  \vspace{2cm}
  \begin{tikzpicture}
    \draw[blue!40!black, thick] (0,0) -- (12,0);
    \draw[blue!40!black, very thick] (0,0.05) -- (12,0.05);
  \end{tikzpicture}
  \vspace{2cm}

  {\Large\bfseries Iñaki Castillo\par}
  \vspace{0.5cm}
  {\large Junio de 2026\par}
  \vfill
  \begin{tikzpicture}
    \draw[blue!40!black, thick] (0,0) -- (12,0);
    \draw[blue!40!black, very thick] (0,-0.05) -- (12,-0.05);
  \end{tikzpicture}
\end{titlepage}

% ── Índice ────────────────────────────────────────────────
\tableofcontents
\newpage

% ─────────────────────────────────────────────────────────
%   INTRODUCCIÓN GENERAL
% ─────────────────────────────────────────────────────────
\section{Introducción}

Este documento recopila y documenta cuatro notebooks de Python orientados al
análisis de series de tiempo. Los temas cubiertos van desde los fundamentos
teóricos —como el proceso de caminata aleatoria— hasta técnicas de modelado
más avanzadas como SARIMA, SARIMAX y Prophet, pasando por el análisis de
series financieras reales obtenidas con \texttt{yfinance}.

La notación general que se utiliza a lo largo del documento es la siguiente:

\begin{center}
\renewcommand{\arraystretch}{1.4}
\begin{tabular}{cl}
  \toprule
  \textbf{Símbolo} & \textbf{Significado} \\
  \midrule
  $\{X_t\}$ & Proceso estocástico (serie de tiempo) \\
  $Z_t$     & Ruido blanco gaussiano: $Z_t \sim \mathcal{N}(0,\sigma^2)$ \\
  $\Delta$  & Operador de diferencia: $\Delta X_t = X_t - X_{t-1}$ \\
  $B$       & Operador de retroceso (\textit{backshift}): $B X_t = X_{t-1}$ \\
  $p, d, q$ & Órdenes AR, integración y MA (no estacional) \\
  $P, D, Q$ & Órdenes AR, integración y MA estacional \\
  $m$       & Periodo estacional \\
  \bottomrule
\end{tabular}
\end{center}

\newpage

% ─────────────────────────────────────────────────────────
%   NOTEBOOK 1: CAMINATA ALEATORIA
% ─────────────────────────────────────────────────────────
\section{Caminata Aleatoria (\textit{Random Walk})}

\subsection{Descripción del Notebook}

El notebook \texttt{caminata\_aleat.ipynb} introduce el concepto de
\textbf{caminata aleatoria}, uno de los procesos estocásticos más importantes
en el análisis de series de tiempo y en finanzas. Se simula el proceso, se
diferencia para obtener una serie estacionaria y se aplica el test de
Dickey-Fuller Aumentado (ADF) para confirmar la estacionariedad.

\subsection{Fundamentos Teóricos}

\begin{definicion}
  Una \textbf{caminata aleatoria} es un proceso estocástico $\{X_t\}$ definido
  por la recurrencia:
  \[
    X_t = X_{t-1} + Z_t, \qquad t = 1, 2, \ldots, T
  \]
  donde $Z_t \overset{iid}{\sim} \mathcal{N}(0,1)$ es ruido blanco gaussiano
  y $X_0$ es un valor inicial fijo (generalmente $X_0 = 0$).
\end{definicion}

La solución explícita del proceso es:
\[
  X_t = X_0 + \sum_{k=1}^{t} Z_k
\]

Sus propiedades estadísticas son:
\begin{align*}
  \mathbb{E}[X_t] &= X_0 = 0 \\
  \mathrm{Var}[X_t] &= t\,\sigma^2
\end{align*}

Como la varianza crece linealmente con el tiempo, la caminata aleatoria
\textbf{no es estacionaria en covarianza}. Esto tiene implicaciones directas
para el modelado: aplicar modelos ARMA directamente a este proceso sería
incorrecto.

\subsubsection{Diferenciación}

Para convertir el proceso en estacionario se aplica el \textbf{operador de
diferencia}:
\[
  \Delta X_t = X_t - X_{t-1} = Z_t
\]

La serie diferenciada $\{\Delta X_t\}$ es ruido blanco con media cero y
varianza constante $\sigma^2$, es decir, \textit{sí} es estacionaria.

\subsubsection{Test de Dickey-Fuller Aumentado (ADF)}

El ADF contrasta las hipótesis:
\[
  H_0: \text{la serie tiene raíz unitaria (no es estacionaria)} \quad
  H_1: \text{la serie es estacionaria}
\]

Se rechaza $H_0$ cuando el estadístico ADF es \textbf{más negativo} que el
valor crítico, o equivalentemente cuando el \textit{p-valor} $< 0.05$.

\subsection{Código: Simulación}

\begin{lstlisting}[caption={Simulación de la caminata aleatoria}]
import numpy as np
import matplotlib.pyplot as plt

# Parámetros
np.random.seed(42)   # Reproducibilidad
T  = 200             # Número de periodos
X0 = 0               # Valor inicial

# Generar ruido blanco Z_t ~ N(0,1)
Z = np.random.normal(0, 1, T)

# Generar proceso X_t = X_{t-1} + Z_t
X = np.zeros(T)
X[0] = X0

for t in range(1, T):
    X[t] = X[t-1] + Z[t]

# Gráfica
plt.figure(figsize=(12, 6), dpi=120)
plt.plot(X, linewidth=2.5)
plt.scatter(range(T), X, s=15)
plt.title("Simulación de Random Walk\n$X_t = X_{t-1} + Z_t$", fontsize=16)
plt.xlabel("Tiempo (t)", fontsize=12)
plt.ylabel("Valor de $X_t$", fontsize=12)
plt.grid(alpha=0.3)
plt.axhline(0, linestyle='--')
plt.tight_layout()
plt.show()
\end{lstlisting}

\subsection{Código: Diferenciación y Test ADF}

\begin{lstlisting}[caption={Diferenciación y Test ADF sobre la caminata aleatoria}]
from statsmodels.tsa.stattools import adfuller

# Diferenciación: Delta X_t = X_t - X_{t-1}
X_diff = np.diff(X)

# Gráfica de la serie diferenciada
plt.figure(figsize=(12, 6), dpi=120)
plt.plot(X_diff, linewidth=2.5)
plt.scatter(range(len(X_diff)), X_diff, s=15)
plt.title("Serie Diferenciada\n$\\Delta X_t = X_t - X_{t-1}$", fontsize=16)
plt.xlabel("Tiempo (t)", fontsize=12)
plt.ylabel("Valor Diferenciado", fontsize=12)
plt.grid(alpha=0.3)
plt.axhline(0, linestyle='--')
plt.tight_layout()
plt.show()

# Test ADF
adf_result = adfuller(X_diff)

print(f"Estadístico ADF : {adf_result[0]:.4f}")
print(f"p-valor         : {adf_result[1]:.4f}")
print("Valores críticos:")
for key, value in adf_result[4].items():
    print(f"   {key}: {value:.4f}")

if adf_result[1] < 0.05:
    print("\nConclusión: La serie diferenciada ES estacionaria (rechazamos H0).")
else:
    print("\nConclusión: La serie diferenciada NO es estacionaria.")
\end{lstlisting}

\begin{resultado}
  La serie $\Delta X_t = Z_t$ es ruido blanco, por lo que el test ADF rechaza
  la hipótesis nula de raíz unitaria. Se espera un estadístico muy negativo y
  un $p$-valor $\ll 0.05$, confirmando que diferenciar \textbf{una sola vez}
  es suficiente para hacer estacionaria a la caminata aleatoria.
\end{resultado}

\newpage

% ─────────────────────────────────────────────────────────
%   NOTEBOOK 2: SERIES FINANCIERAS
% ─────────────────────────────────────────────────────────
\section{Series Financieras — Amazon (AMZN)}

\subsection{Descripción del Notebook}

El notebook \texttt{Series\_Financieras.ipynb} descarga precios históricos de
la acción de Amazon (ticker: \texttt{AMZN}) mediante la librería
\texttt{yfinance}, visualiza la serie de precios de cierre con \texttt{Altair}
y calcula los \textbf{log-retornos} como transformación estándar en finanzas
cuantitativas.

\subsection{Fundamentos Teóricos}

\subsubsection{Precios de cierre ajustados}

En finanzas cuantitativas se trabaja habitualmente con el precio de cierre
ajustado $P_t$ que incorpora dividendos y \textit{splits}. El precio crudo es
no estacionario (generalmente con tendencia estocástica), por lo que no puede
emplearse directamente en modelos ARMA sin transformación.

\subsubsection{Log-retornos}

Sea $P_t$ el precio de cierre en el periodo $t$. El \textbf{log-retorno
continuo} se define como:
\[
  r_t = \ln\!\left(\frac{P_t}{P_{t-1}}\right) = \ln P_t - \ln P_{t-1}
\]

Propiedades importantes de los log-retornos:
\begin{itemize}[noitemsep]
  \item Son \textbf{aditivos en el tiempo}: el retorno acumulado en $[1,T]$ es
        $\sum_{t=1}^{T} r_t$.
  \item Para retornos pequeños, $r_t \approx (P_t - P_{t-1})/P_{t-1}$ (retorno
        aritmético).
  \item Empíricamente presentan \textit{colas pesadas} y
        \textit{agrupamiento de volatilidad} (\textit{volatility clustering}).
  \item Son aproximadamente estacionarios en media (aunque su varianza puede
        variar, capturada por modelos ARCH/GARCH).
\end{itemize}

\subsection{Código}

\begin{lstlisting}[caption={Descarga de datos y cálculo de log-retornos de AMZN}]
import yfinance as yf
import pandas as pd
import numpy as np
import altair as alt

# ── Descarga de datos ──────────────────────────────────
amzn = yf.download("AMZN", start="2010-01-01")
amzn_precios = amzn["Close"]

# ── Visualización interactiva con Altair ───────────────
df = amzn_precios.reset_index()
df.columns = ["Date", "Close"]

chart = (
    alt.Chart(df)
    .mark_line()
    .encode(
        x=alt.X("Date:T", title="Fecha"),
        y=alt.Y("Close:Q", title="Precio de cierre (USD)"),
        tooltip=["Date:T", "Close:Q"],
    )
    .properties(title="Serie de tiempo AMZN", width=900, height=400)
    .interactive()
)
chart

# ── Log-retornos ──────────────────────────────────────
# r_t = ln(P_t) - ln(P_{t-1})
amzn_log_returns = np.log(amzn_precios / amzn_precios.shift(1))
amzn_log_returns = amzn_log_returns.dropna()   # Eliminar primer NaN

display(amzn_log_returns.head())
\end{lstlisting}

\begin{nota}
  El primer valor del vector de log-retornos resulta \texttt{NaN} porque no
  existe $P_{t-1}$ para $t=0$. Se elimina con \texttt{.dropna()} antes de
  cualquier análisis posterior. En la práctica, los log-retornos de AMZN desde
  2010 presentan el comportamiento típico de renta variable: media cercana a
  cero, exceso de curtosis (colas pesadas) y períodos de alta y baja
  volatilidad.
\end{nota}

\newpage

% ─────────────────────────────────────────────────────────
%   NOTEBOOK 3: ESTACIONALIDAD Y SARIMAX
% ─────────────────────────────────────────────────────────
\section{Estacionalidad y SARIMAX}

\subsection{Descripción del Notebook}

El notebook \texttt{Seasonality.ipynb} construye un dataset sintético con
\textbf{tendencia}, \textbf{estacionalidad anual} y una \textbf{variable
exógena} (efecto de promociones) para motivar el uso del modelo SARIMAX, que
es la extensión de SARIMA con variables explicativas externas.

\subsection{Fundamentos Teóricos}

\subsubsection{Descomposición de una serie temporal}

Una serie de tiempo $\{Y_t\}$ puede descomponerse como:
\[
  Y_t = T_t + S_t + \varepsilon_t
\]
donde $T_t$ es la \textbf{tendencia}, $S_t$ es el componente
\textbf{estacional} y $\varepsilon_t$ es el término de \textbf{error}.

\subsubsection{Modelo SARIMAX}

El modelo SARIMAX$\,(p,d,q)(P,D,Q)_m$ extiende SARIMA incluyendo regresores
exógenos $\mathbf{x}_t$:
\[
  \Phi_P(B^m)\,\phi_p(B)\,\Delta^d\Delta_m^D Y_t
  = \Theta_Q(B^m)\,\theta_q(B)\,\varepsilon_t + \boldsymbol{\beta}^\top \mathbf{x}_t
\]

donde:
\begin{align*}
  \phi_p(B)    &= 1 - \phi_1 B - \cdots - \phi_p B^p & &\text{(polinomio AR no estacional)} \\
  \theta_q(B)  &= 1 + \theta_1 B + \cdots + \theta_q B^q & &\text{(polinomio MA no estacional)} \\
  \Phi_P(B^m)  &= 1 - \Phi_1 B^m - \cdots - \Phi_P B^{Pm} & &\text{(polinomio AR estacional)} \\
  \Theta_Q(B^m)&= 1 + \Theta_1 B^m + \cdots + \Theta_Q B^{Qm} & &\text{(polinomio MA estacional)} \\
  \Delta^d     &= (1-B)^d & &\text{(diferenciación no estacional)} \\
  \Delta_m^D   &= (1-B^m)^D & &\text{(diferenciación estacional)}
\end{align*}

\subsection{Generación del Dataset Sintético}

\begin{lstlisting}[caption={Construcción del dataset sintético con tendencia, estacionalidad y variables exógenas}]
import pandas as pd
import numpy as np
import matplotlib.pyplot as plt

# Rango de fechas: 5 años, datos mensuales
dates = pd.date_range(start='2018-01-01', periods=5*12, freq='MS')
np.random.seed(42)

# Componentes de la serie
trend        = np.linspace(0, 100, len(dates))                              # Tendencia lineal
seasonality  = (15 * np.sin(np.linspace(0, 3*2*np.pi, len(dates)))          # Ciclo anual
              + 10 * np.cos(np.linspace(0, 2*np.pi, len(dates))))
noise        = np.random.normal(0, 5, len(dates))

# Variable exógena: efecto de promociones
promotion_effect = np.zeros(len(dates))
promotion_months = [3, 8, 15, 20, 27, 32, 39, 44, 51, 56]

for month_idx in promotion_months:
    if month_idx < len(dates):
        promotion_effect[month_idx] = 20 + np.random.normal(0, 3)

# Serie final: combinación aditiva
time_series = trend + seasonality + promotion_effect + noise

# DataFrame
df = pd.DataFrame(
    {'value': time_series, 'promotion': promotion_effect > 0},
    index=dates
)
display(df.head())

# Visualización
plt.figure(figsize=(15, 7))
plt.plot(df.index, df['value'], label='Serie Temporal Sintética')
plt.scatter(df.index[df['promotion']], df['value'][df['promotion']],
            color='red', s=50, label='Mes con Promoción')
plt.title('Serie Temporal Sintética con Tendencia, Estacionalidad y Promociones')
plt.xlabel('Fecha')
plt.ylabel('Valor')
plt.legend()
plt.grid(True)
plt.show()
\end{lstlisting}

\begin{nota}
  En escenarios reales, la columna \texttt{promotion} puede reemplazarse por
  cualquier variable exógena relevante: días festivos, índices económicos,
  temperatura, campañas de marketing, etc. SARIMAX aprovecha esta información
  para mejorar la precisión del pronóstico en comparación con un SARIMA puro.
\end{nota}

\newpage

% ─────────────────────────────────────────────────────────
%   NOTEBOOK 4: SARIMA Y PROPHET
% ─────────────────────────────────────────────────────────
\section{Modelo SARIMA y Prophet}

\subsection{Descripción del Notebook}

El notebook \texttt{Clase\_1.ipynb} implementa paso a paso:
\begin{enumerate}[noitemsep]
  \item Generación de una serie de tiempo sintética con tendencia y estacionalidad.
  \item Ajuste de un modelo \textbf{SARIMA}$(0,1,1)(0,1,1)_{12}$ (Airline Model).
  \item Pronóstico a 24 meses.
  \item Análisis de residuos (ACF, PACF, Q-Q plot, tests estadísticos).
  \item Pronóstico alternativo con el modelo \textbf{Prophet} de Meta.
\end{enumerate}

\subsection{Fundamentos Teóricos}

\subsubsection{Modelo SARIMA$(0,1,1)(0,1,1)_{12}$}

Conocido como el \textit{Airline Model} (Box \& Jenkins, 1976), es uno de los
modelos de series de tiempo más utilizados en la práctica. Su ecuación es:

\[
  (1-B)(1-B^{12})\,Y_t
  = (1 + \theta_1 B)(1 + \Theta_1 B^{12})\,\varepsilon_t
\]

Las componentes de este modelo específico son:

\begin{center}
\renewcommand{\arraystretch}{1.4}
\begin{tabular}{cll}
  \toprule
  \textbf{Parámetro} & \textbf{Valor} & \textbf{Interpretación} \\
  \midrule
  $p$ & 0 & Sin AR no estacional \\
  $d$ & 1 & Una diferenciación no estacional \\
  $q$ & 1 & MA de orden 1 no estacional \\
  $P$ & 0 & Sin AR estacional \\
  $D$ & 1 & Una diferenciación estacional \\
  $Q$ & 1 & MA de orden 1 estacional \\
  $m$ & 12 & Periodo anual (datos mensuales) \\
  \bottomrule
\end{tabular}
\end{center}

El doble operador de diferencia $\Delta\Delta_{12}$ elimina simultáneamente
la \textbf{tendencia} y la \textbf{estacionalidad} de la serie.

\subsubsection{Diagnóstico de Residuos}

Para validar el ajuste del modelo los residuos $\hat{\varepsilon}_t$ deben
comportarse como \textbf{ruido blanco}:

\begin{itemize}[noitemsep]
  \item \textbf{Ljung-Box}: contrasta $H_0$: no hay autocorrelación en los
        residuos. $p > 0.05$ es deseable.
  \item \textbf{Jarque-Bera}: contrasta $H_0$: los residuos son normales.
        $p > 0.05$ es deseable.
  \item \textbf{Heteroscedasticidad}: contrasta $H_0$: varianza constante.
        $p > 0.05$ es deseable.
\end{itemize}

\subsubsection{Modelo Prophet}

Prophet es un modelo de pronóstico aditivo desarrollado por Meta (2017):
\[
  Y(t) = g(t) + s(t) + h(t) + \varepsilon_t
\]

donde $g(t)$ es la tendencia, $s(t)$ los componentes estacionales (Fourier),
$h(t)$ el efecto de días festivos y $\varepsilon_t$ el error. Es especialmente
robusto ante datos faltantes y cambios de tendencia.

\subsection{Código: Generación de Datos y Ajuste SARIMA}

\begin{lstlisting}[caption={Generación de la serie sintética}]
import pandas as pd
import numpy as np
import matplotlib.pyplot as plt

np.random.seed(42)
n_points = 120   # 10 años de datos mensuales
index = pd.date_range(start='2010-01-01', periods=n_points, freq='MS')

# Componentes
trend      = np.linspace(0, 20, n_points)
seasonal   = 10 * np.sin(np.linspace(0, 3 * np.pi, n_points) * 2)
noise      = np.random.normal(0, 1.5, n_points)

# Serie resultante
data        = trend + seasonal + noise + 50
time_series = pd.Series(data, index=index)

plt.figure(figsize=(12, 6))
plt.plot(time_series)
plt.title('Serie de Tiempo Sintética')
plt.xlabel('Fecha')
plt.ylabel('Valor')
plt.grid(True)
plt.show()
\end{lstlisting}

\begin{lstlisting}[caption={Ajuste del modelo SARIMA$(0,1,1)(0,1,1)_{12}$}]
from statsmodels.tsa.statespace.sarimax import SARIMAX

order          = (0, 1, 1)
seasonal_order = (0, 1, 1, 12)

model     = SARIMAX(time_series, order=order, seasonal_order=seasonal_order,
                    enforce_stationarity=False, enforce_invertibility=False)
model_fit = model.fit(disp=False)

print(model_fit.summary())
\end{lstlisting}

\subsection{Código: Pronóstico SARIMA}

\begin{lstlisting}[caption={Pronóstico a 24 meses con SARIMA}]
n_forecast = 24

forecast = model_fit.predict(
    start=len(time_series),
    end=len(time_series) + n_forecast - 1
)

forecast_index = pd.date_range(
    start=time_series.index[-1] + pd.DateOffset(months=1),
    periods=n_forecast, freq='MS'
)
forecast.index = forecast_index

plt.figure(figsize=(14, 7))
plt.plot(time_series, label='Datos Originales')
plt.plot(forecast, label='Pronóstico SARIMA', color='red', linestyle='--')
plt.title('Serie Sintética con Pronóstico SARIMA')
plt.xlabel('Fecha')
plt.ylabel('Valor')
plt.legend()
plt.grid(True)
plt.show()
\end{lstlisting}

\subsection{Código: Análisis de Residuos}

\begin{lstlisting}[caption={Diagnóstico gráfico y estadístico de los residuos}]
import seaborn as sns
import statsmodels.api as sm
from statsmodels.graphics.tsaplots import plot_acf, plot_pacf

residuals = model_fit.resid

# ── Gráficas de diagnóstico ────────────────────────────
fig, axes = plt.subplots(3, 1, figsize=(16, 10))

axes[0].plot(residuals)
axes[0].set_title('Residuos a lo Largo del Tiempo')
axes[0].grid(True)

sns.histplot(residuals, kde=True, ax=axes[1])
axes[1].set_title('Distribución de los Residuos')

sm.qqplot(residuals, line='s', ax=axes[2])
axes[2].set_title('Q-Q Plot de los Residuos')

plt.tight_layout()
plt.show()

# ── ACF y PACF ────────────────────────────────────────
fig2, (ax1, ax2) = plt.subplots(2, 1, figsize=(14, 8))
plot_acf(residuals,  lags=40, ax=ax1, zero=False)
plot_pacf(residuals, lags=40, ax=ax2, zero=False)
plt.tight_layout()
plt.show()
\end{lstlisting}

\begin{resultado}
  Para el modelo ajustado se esperan los siguientes valores en los tests del
  \texttt{summary()}: Ljung-Box $p \approx 0.76$, Jarque-Bera $p \approx 0.53$,
  Heteroscedasticidad $p \approx 0.64$. Todos son $> 0.05$, lo que indica que
  los residuos se comportan como ruido blanco y el modelo es adecuado.
\end{resultado}

\subsection{Código: Pronóstico con Prophet}

\begin{lstlisting}[caption={Ajuste y pronóstico con el modelo Prophet}]
from prophet import Prophet

# Prophet requiere columnas 'ds' (fecha) e 'y' (valor)
prophet_df = time_series.reset_index()
prophet_df.columns = ['ds', 'y']

# Ajuste del modelo
model_prophet = Prophet(yearly_seasonality=True)
model_prophet.fit(prophet_df)

# DataFrame de fechas futuras
future = model_prophet.make_future_dataframe(periods=n_forecast, freq='MS')
forecast_prophet = model_prophet.predict(future)

# Columnas principales del pronóstico
display(forecast_prophet[['ds', 'yhat', 'yhat_lower', 'yhat_upper']].tail())

# Visualización
fig = model_prophet.plot(forecast_prophet)
plt.title('Serie Sintética con Pronóstico de Prophet')
plt.xlabel('Fecha')
plt.ylabel('Valor')
plt.show()

# Descomposición de componentes
fig2 = model_prophet.plot_components(forecast_prophet)
plt.show()
\end{lstlisting}

\begin{nota}
  Prophet descompone automáticamente el pronóstico en tendencia y
  estacionalidad. Las columnas \texttt{yhat\_lower} y \texttt{yhat\_upper}
  representan el intervalo de credibilidad al 80\,\% por defecto. A diferencia
  de SARIMA, Prophet no requiere que la serie sea estacionaria ni que el
  usuario especifique $p$, $d$, $q$.
\end{nota}

\newpage

% ─────────────────────────────────────────────────────────
%   COMPARACIÓN DE MODELOS
% ─────────────────────────────────────────────────────────
\section{Comparación de Modelos}

\begin{center}
\renewcommand{\arraystretch}{1.5}
\begin{tabular}{lcccc}
  \toprule
  \textbf{Característica} & \textbf{SARIMA} & \textbf{SARIMAX} & \textbf{Prophet} \\
  \midrule
  Variables exógenas       & No  & Sí  & Días festivos \\
  Estacionalidad múltiple  & No  & No  & Sí \\
  Requiere estacionariedad & Sí  & Sí  & No \\
  Selección de hiperparámetros & Manual & Manual & Automática \\
  Intervalo de predicción  & Frecuentista & Frecuentista & Bayesiano \\
  Robusto a datos faltantes & No & No & Sí \\
  \bottomrule
\end{tabular}
\end{center}

\newpage

% ─────────────────────────────────────────────────────────
%   DEPENDENCIAS
% ─────────────────────────────────────────────────────────
\section{Dependencias y Entorno}

Las siguientes librerías de Python son necesarias para reproducir todos los
notebooks:

\begin{lstlisting}[language=bash, caption={Instalación de dependencias}]
pip install numpy pandas matplotlib seaborn statsmodels yfinance altair prophet
\end{lstlisting}

\begin{center}
\renewcommand{\arraystretch}{1.4}
\begin{tabular}{lll}
  \toprule
  \textbf{Librería} & \textbf{Uso principal} & \textbf{Versión recomendada} \\
  \midrule
  \texttt{numpy}       & Álgebra lineal y generación de datos & $\geq 1.24$ \\
  \texttt{pandas}      & Manejo de series temporales & $\geq 2.0$ \\
  \texttt{matplotlib}  & Visualización estática & $\geq 3.7$ \\
  \texttt{seaborn}     & Visualización estadística & $\geq 0.12$ \\
  \texttt{statsmodels} & SARIMA, ADF, ACF, PACF & $\geq 0.14$ \\
  \texttt{yfinance}    & Descarga de datos financieros & $\geq 0.2$ \\
  \texttt{altair}      & Gráficas interactivas & $\geq 5.0$ \\
  \texttt{prophet}     & Modelo Prophet (Meta) & $\geq 1.1$ \\
  \bottomrule
\end{tabular}
\end{center}

% ─────────────────────────────────────────────────────────
%   REFERENCIAS
% ─────────────────────────────────────────────────────────
\section*{Referencias}
\addcontentsline{toc}{section}{Referencias}

\begin{enumerate}[label={[\arabic*]}]
  \item Box, G.E.P., Jenkins, G.M., Reinsel, G.C. \& Ljung, G.M. (2015).
        \textit{Time Series Analysis: Forecasting and Control}. 5th ed. Wiley.
  \item Hyndman, R.J. \& Athanasopoulos, G. (2021).
        \textit{Forecasting: Principles and Practice}. 3rd ed. OTexts.
        \url{https://otexts.com/fpp3/}
  \item Taylor, S.J. \& Letham, B. (2018). Forecasting at Scale.
        \textit{The American Statistician}, 72(1), 37--45.
  \item Seabold, S. \& Perktold, J. (2010). Statsmodels: Econometric and
        Statistical Modeling with Python. \textit{9th Python in Science Conf.}
\end{enumerate}

\end{document}
